\documentclass[11pt]{article}

\usepackage[margin=1in]{geometry}
\usepackage{amsmath,amssymb}
\usepackage{booktabs}
\usepackage{array}
\usepackage{enumitem}
\usepackage{xcolor}
\PassOptionsToPackage{hyphens}{url}
\usepackage{hyperref}

\hypersetup{
  colorlinks=true,
  linkcolor=blue!60!black,
  citecolor=blue!60!black,
  urlcolor=blue!60!black
}

\newcommand{\R}{\mathbb{R}}
\newcommand{\Vactive}{V_{\mathrm{act}}}
\newcommand{\frontier}{F}
\newcommand{\dens}{\operatorname{dens}}
\newcommand{\argmax}{\operatorname*{arg\,max}}

\title{Densest Community Search in Directed Citation Graphs via Lazy Branch-and-Price}
\author{Project Draft}
\date{\today}

\begin{document}
\maketitle

\begin{abstract}
Citation graphs are a natural substrate for exploratory scholarly search:
given a paper, a user often wants a compact set of nearby papers that
forms a coherent research community. Simple local expansion is cheap, but
it does not optimize cohesion, satisfy a target size, or guarantee a
formal connectivity property. This paper presents a query-anchored
densest community search system for directed citation graphs. Given a
directed graph, a query paper, and a minimum community size $k \ge 3$,
the system seeks a set containing the query that maximizes induced
directed edge density $|E(S)| / (|S|(|S|-1))$. The solver also supports
a rooted $\kappa$-edge-connectivity constraint.

The method combines Dinkelbach's fractional-programming transform with a
lazy branch-and-price algorithm. Each Dinkelbach iteration solves a
parametric integer program over a dynamically expanded active set. The
restricted master problem uses node, directed-edge, and unordered-pair
variables; column generation promotes promising frontier nodes;
branch-and-bound enforces integrality; Boolean quadratic polytope
triangle cuts tighten fractional relaxations; and min-cut separation
enforces optional rooted connectivity. The implementation supports both
local CSV simulation and live OpenAlex citation retrieval, wrapped by a
FastAPI/React interface that streams optimization telemetry.

Evaluation has three perspectives on Cora-ML (CitationFull). The
cluster-quality perspective, on the $804$ val+test queries, has
branch-and-price (BP) returning the densest and most algebraically
connected compact communities at every $k \in \{3, 4, 5\}$: at $k = 3$,
BP at $\kappa = 0$ reaches edge density $0.51 \pm 0.10$ and BP at
$\kappa = 2$ reaches $\lambda_2 = 1.27 \pm 0.27$, against BFS depth-$1$
at density $0.35 \pm 0.11$ and AvgDeg at density $0.21 \pm 0.12$ with
size $21.2 \pm 15.2$ and average degree $3.02 \pm 1.55$. The cost
perspective has BFS sub-second, AvgDeg under $0.1$ s, BP heterogeneous
with mean wall time growing from $1.5$ s at $k = 3, \kappa = 0$ to
$96$ s at $k = 5, \kappa = 2$ under a $300$ s hard cap. The
classification perspective tunes on the $402$ val queries and evaluates
on the $402$ test queries: BFS depth-$1$ at $k = 4$ wins macro F1
($0.856 \pm 0.019$), BP at $k = 4, \kappa = 2$ runs second
($0.811 \pm 0.021$), AvgDeg at $k = 4$ trails ($0.603 \pm 0.024$). On
the $59$-node BFS-hard test subset, where the depth-$1$ vote disagrees
with the true label by construction, BP leads macro F1
($0.190 \pm 0.056$) ahead of AvgDeg ($0.146 \pm 0.056$) and BFS
($0.080 \pm 0.042$). The pattern supports a two-regime reading: direct
neighbours dominate easy citation classification, while density-
constrained optimization is the principled retrieval primitive when
size, density, and connectivity are part of the task.
\end{abstract}

\section{Motivation and Problem Statements}

Citation graphs encode explicit intellectual relationships among papers. A
paper's outgoing citations represent prior work it builds on, while incoming
citations represent later papers that build on it. Many scholarly interfaces
surface these relationships as flat lists: references, citing papers, related
papers, or high-citation neighbors. These views are useful, but they leave open
a graph optimization question: for a given query paper, which compact set of
papers forms the densest citation community around it?

The simplest answer is to return a one-hop neighborhood. In the Cora-ML
experiments reported here, this baseline is strong for classification
because test papers directly cite labeled training papers. But
one-hop expansion is a
heuristic, not an optimization problem. It may return too few papers for a
desired reading set, it may fail a formal size constraint, it may return a set
whose internal edge density is far from optimal, and it gives no way to trade
off compactness, density, and connectivity.

This project studies a query-anchored densest community search problem for
directed citation graphs. Given a directed graph $G=(V,E)$, a query node $q$,
and a minimum size $k$, we seek a set $S \subseteq V$ such that $q \in S$,
$|S| \geq k$, and the induced directed density
\[
  \dens(S) = \frac{|E(S)|}{|S|(|S|-1)}
\]
is maximized. An optional $\kappa$ parameter additionally requires selected
nodes to have rooted edge connectivity from the query in the undirected support
of the selected citation edges.

The implementation is designed for two settings. In offline simulation mode, the
graph is a local CSV edge list. In live exploration mode, the graph is implicit:
the solver queries OpenAlex for the predecessor and successor papers it needs.
This motivates a lazy algorithm. Rather than loading the whole graph, the solver
starts from a query-centered active set and only promotes frontier nodes that
have positive reduced cost in the current LP relaxation.

The main contributions are:
\begin{enumerate}[leftmargin=1.2cm]
  \item A directed, query-anchored, at-least-$k$ dense-community formulation
  with optional rooted $\kappa$-edge-connectivity.
  \item A lazy branch-and-price algorithm that combines Dinkelbach iterations,
  restricted master LPs, frontier pricing, branch-and-bound, triangle-cut
  separation, incumbent pruning, and min-cut connectivity cuts.
  \item A reproducible node-classification and cluster-quality evaluation
  pipeline on CitationFull datasets, instantiated here on Cora-ML with BFS and
  exact average-degree baselines.
  \item A working interactive system with live OpenAlex citation retrieval and
  streaming solver telemetry.
\end{enumerate}

The empirical claim is intentionally measured. The current Cora-ML evidence does
not show that branch-and-price replaces BFS as a universal classifier. Rather,
BFS and dense-community search solve different problems. BFS is excellent when
direct citation labels are sufficient. Branch-and-price is useful when the
desired object is a compact, dense, query-containing community satisfying
explicit structural constraints.

\subsection{Related Work}

\paragraph{Dense subgraph discovery.}
The unconstrained maximum density subgraph problem is classical. Goldberg gave a
flow-based algorithm for finding a maximum density subgraph
\cite{goldberg1984}. Charikar later developed greedy approximation algorithms
for finding dense graph components \cite{charikar2000}. The size-constrained
densest $k$-subgraph problem is harder and remains a central object in dense
subgraph research. Exact and approximation approaches for dense $k$-subgraph
variants continue to be studied, especially for large graphs
\cite{veremyev2017}.

The present project differs from standard densest subgraph formulations in three
ways. First, it is query anchored: the selected community must contain a
particular paper. Second, it is directed: a citation from $u$ to $v$ is not the
same as a citation from $v$ to $u$. Third, the graph may be accessed lazily
through an API, so the algorithm must reason about which candidate nodes to
materialize.

\paragraph{Fractional programming and branch-and-price.}
The density objective is a ratio. Dinkelbach's method transforms a nonlinear
fractional program into a sequence of parametric optimization problems
\cite{dinkelbach1967}. This is a natural fit for dense subgraph objectives:
instead of optimizing $|E(S)|/(|S|(|S|-1))$ directly, each iteration optimizes
$|E(S)|-\lambda |S|(|S|-1)$ for a fixed density estimate $\lambda$.

Branch-and-price combines branch-and-bound with column generation and is used
for integer programs whose full variable set is too large to enumerate at once
\cite{barnhart1998}. In this project, the variable universe is large not only
because the graph may be large, but also because live OpenAlex mode discovers
nodes on demand. Column generation therefore acts as both an optimization device
and a graph-expansion policy.

\paragraph{Citation data.}
The live mode uses OpenAlex, a fully open scholarly index and API
\cite{priem2022}. The Works API exposes citation relationships through outgoing
references and incoming cited-by queries \cite{openalexworks}. The offline
classification evaluation uses PyTorch Geometric's \texttt{CitationFull}
dataset loader \cite{fey2019,pygcitationfull}, which provides citation datasets
from the Graph2Gauss benchmark collection \cite{bojchevski2017}.

\subsection{Problem Statement}

Let $G=(V,E)$ be a directed graph. In the citation setting, $(u,v)\in E$ means
that paper $u$ cites paper $v$. Let $q\in V$ be the query paper and $k\geq 3$
be a minimum target size. For $S\subseteq V$, define the induced directed edge
set
\[
  E(S) = \{(u,v)\in E : u\in S,\ v\in S,\ u\neq v\}.
\]
The directed density of $S$ is
\[
  \dens(S) = \frac{|E(S)|}{|S|(|S|-1)}.
\]
The primary optimization problem is
\begin{align}
  \max_{S\subseteq V} \quad & \dens(S) \label{eq:main-problem}\\
  \text{s.t.} \quad & q \in S, \nonumber\\
                    & |S| \geq k. \nonumber
\end{align}

The use of $|S|\geq k$ rather than $|S|=k$ is important. A scholarly community
may naturally contain more than $k$ papers if those papers preserve density. The
solver therefore optimizes over a minimum-size family and then prunes
unproductive extra nodes from integer incumbents.

For $\kappa>0$, the solver adds a rooted connectivity requirement. Let
$H(S)$ be the undirected support graph obtained from selected directed internal
edges. The connectivity constraint requires every $v\in S\setminus\{q\}$ to
have at least $\kappa$ edge-disjoint paths from $q$ in $H(S)$. By Menger's
theorem, this can be checked through min-cut/max-flow computations and enforced
by adding cut-set inequalities when a candidate violates the condition.

\section{Approaches}

\subsection{Oracle-Based Lazy Graph Access}

The solver is organized around an \texttt{IGraphOracle}. Given a node, the
oracle returns predecessor and successor lists. In simulation mode, these lists
come from a local CSV edge list. In OpenAlex mode, they are fetched from the
OpenAlex API. Incoming retrieval is capped by \texttt{max\_in\_edges}; setting
this value to zero disables incoming expansion. The branch-and-price algorithm
never assumes the full graph is preloaded.

Internally the solver maintains:
\begin{itemize}[leftmargin=1.1cm]
  \item an active set $\Vactive$ of nodes currently represented in the restricted
  master problem,
  \item a frontier $\frontier$ of adjacent nodes discovered but not yet promoted,
  \item cached outgoing and incoming adjacency maps for queried nodes,
  \item a list of pending directed edges whose endpoint variables may not both
  exist yet.
\end{itemize}

Initialization performs a BFS from $q$ until at least $k$ active nodes are
available. It then queries active nodes that were discovered as neighbors but
not expanded during the first BFS pass, so their adjacency lists are completed.
Nodes that fail oracle retrieval are blacklisted. If the query node itself
cannot be fetched, the solve aborts.

This active-set design is the practical bridge between optimization and live
retrieval. The solver can begin with a small local neighborhood, solve a
restricted relaxation, and then ask which frontier papers are worth the next API
queries.

\subsection{Dinkelbach Outer Loop}

The density objective is fractional. For a current density estimate $\lambda$,
define the parametric objective
\begin{equation}
  \phi_\lambda(S) = |E(S)| - \lambda |S|(|S|-1).
  \label{eq:parametric}
\end{equation}
Dinkelbach's method repeatedly solves
\begin{align}
  \max_{S\subseteq V}\quad & \phi_\lambda(S)\\
  \text{s.t.}\quad & q\in S,\quad |S|\geq k,
\end{align}
then updates $\lambda \leftarrow \dens(S)$ if an improving solution is found.
If the best parametric value is nonpositive, no set improves the current density
estimate for the current explored problem.

The implementation initializes $\lambda$ from a query-centered BFS seed after a
greedy pruning pass. This matters because a stronger initial density makes the
first parametric problem tighter and can reduce branch-and-bound work. Each
Dinkelbach iteration calls the branch-and-price solver. The loop stops when the
parametric solve finds no positive improvement, when density no longer
increases beyond tolerance, or when the configured maximum number of iterations
is reached.

\subsection{Restricted Master Formulation}

For a fixed $\lambda$ and active set $\Vactive$, the solver forms a restricted
linearized model. The binary variables are:
\[
  x_v \in \{0,1\}\quad(v\in \Vactive),
\]
where $x_v=1$ means node $v$ is selected;
\[
  y_{uv}\in \{0,1\}\quad((u,v)\in E,\ u,v\in \Vactive),
\]
where $y_{uv}=1$ means the directed edge $(u,v)$ is internal; and
\[
  w_{uv}\in \{0,1\}\quad(\{u,v\}\subseteq \Vactive,\ u<v),
\]
where $w_{uv}=1$ represents $x_u x_v$.

For integer $x$, the parametric objective can be written as
\begin{equation}
  \sum_{(u,v)\in E(\Vactive)} y_{uv}
  - 2\lambda \sum_{\{u,v\}\subseteq \Vactive,\ u<v} w_{uv}.
  \label{eq:rmp-obj}
\end{equation}
The factor of 2 appears because
$|S|(|S|-1)=2\binom{|S|}{2}$.

The restricted master problem is
\begin{align}
  \max \quad &
  \sum_{(u,v)\in E(\Vactive)} y_{uv}
  -2\lambda\sum_{u<v}w_{uv}
  \label{eq:rmp}\\
  \text{s.t.}\quad
  & \sum_{v\in\Vactive} x_v \geq k, \nonumber\\
  & x_q = 1, \nonumber\\
  & y_{uv} \leq x_u,\quad y_{uv}\leq x_v
    && ((u,v)\in E(\Vactive)), \nonumber\\
  & w_{uv} \geq x_u+x_v-1
    && (u<v,\ u,v\in\Vactive), \nonumber\\
  & x_v,y_{uv},w_{uv}\in\{0,1\}. \nonumber
\end{align}

The formulation uses only the linearization constraints needed by the signs of
the objective. Since $y$ has positive coefficient, maximizing sets $y_{uv}=1$
whenever both endpoints are selected, and $y_{uv}=0$ otherwise. Since $w$ has
negative coefficient, maximizing sets $w_{uv}$ as low as allowed, which equals
$1$ when both endpoints are selected and $0$ otherwise for integer $x$.

During column generation and branch-and-bound, the integrality constraints are
relaxed to $[0,1]$. This relaxation is not exact, so the algorithm tightens it
with branching and cuts.

\subsection{Column Generation as Frontier Expansion}

Solving the LP over $\Vactive$ yields fractional node values $\bar{x}$ and a
dual value $\pi$ for the size constraint. The solver then prices each frontier
node $f\in\frontier$. Let
\[
  \operatorname{fracdeg}(f)
  =
  \sum_{u:\ (u,f)\in E,\ u\in\Vactive} \bar{x}_u
  +
  \sum_{v:\ (f,v)\in E,\ v\in\Vactive} \bar{x}_v
\]
be the fractional internal degree that $f$ would gain against the current active
solution. The implemented reduced-cost score is
\begin{equation}
  \operatorname{rc}(f)
  =
  \operatorname{fracdeg}(f)
  -2\lambda\sum_{v\in\Vactive}\bar{x}_v
  -\pi.
  \label{eq:pricing}
\end{equation}
The first term estimates the directed internal-edge gain from adding $f$. The
second term estimates the pairwise density penalty created by pairing $f$ with
currently selected active nodes. The third term accounts for the size constraint
dual. Frontier nodes with positive reduced cost are promoted to $\Vactive$,
queried through the oracle, and synchronized into the restricted master problem.

The implementation promotes nodes in batches rather than one at a time. The
batch size is controlled by a fraction of the active-set size and by minimum and
maximum caps. This reduces model synchronization overhead and API round trips.
Operationally, column generation therefore decides both which LP columns to add
and which citation neighborhoods to fetch next.

\subsection{Branch-and-Bound}

Each Dinkelbach subproblem is solved by a depth-first branch-and-bound tree. A
branch-and-bound node consists of two lists: nodes fixed to one and nodes fixed
to zero. The root fixes $q$ to one. Before each LP solve, the solver applies
these bounds to the restricted master variables.

After column generation reaches a local LP optimum, the LP bound is compared
with the incumbent parametric value. Nodes whose bound cannot improve the
incumbent are pruned. If the LP solution is fractional, the solver selects a
branching variable using a density-aware rule. For each fractional node $v$, it
computes a fractional internal degree into the current LP solution. Nodes with
internal degree below $2\lambda$ are treated as ``hanging'' nodes: including them
is likely to hurt density, so the zero branch is explored first. If no such node
exists, the solver branches on a more core fractional node and explores the one
branch first.

This branching rule is heuristic, but it reflects the parametric objective. A
node must add enough incident internal edge value to compensate for the
$2\lambda$ marginal pair penalty. Nodes below that threshold are natural
candidates for exclusion.

\subsection{Triangle Cuts for the Pair Relaxation}

The pair variables $w_{uv}$ are lower-bound McCormick variables. In fractional
solutions, the relaxation can be loose: the LP may spread mass over several
nodes while keeping pair variables artificially small. To reduce this effect,
the solver separates selected Boolean quadratic polytope triangle inequalities:
\begin{equation}
  x_u+x_v+x_w-w_{uv}-w_{vw}-w_{uw}\leq 1.
  \label{eq:triangle}
\end{equation}
These inequalities are valid for binary $x$ and $w=x_ix_j$. The implementation
checks triples of moderately fractional nodes, caches $w$ values to avoid
repeated solver lookups, and adds at most 20 violated cuts per separation round.
Cuts are attempted only when there are enough fractional nodes and the current
LP gap justifies the extra separation work.

\subsection{Rooted \texorpdfstring{$\kappa$}{kappa}-Connectivity Separation}

When $\kappa=0$, integer LP solutions can be evaluated directly. When
$\kappa>0$, an integer candidate must pass a rooted edge-connectivity check. The
solver builds an undirected unit-capacity support graph from selected internal
directed edges. For every selected target node $t\neq q$, it computes a max flow
from $q$ to $t$. If the flow is less than $\kappa$, the residual graph gives a
violated cut. Let $R$ be the nodes reachable from $q$ in the residual graph.
The solver adds the cut-set inequality
\begin{equation}
  \sum_{\{u,v\}\in\delta(R)} y_{uv} \geq \kappa x_t,
  \label{eq:kappa-cut}
\end{equation}
where $\delta(R)$ is the set of selected-support edges crossing the cut. The
current branch-and-bound node is then reprocessed with the new inequality. This
is a lazy cut strategy: connectivity constraints are generated only when a
candidate solution violates them.

\subsection{Incumbent Pruning}

Whenever the solver obtains a discrete solution, it tries to remove non-query
nodes that improve the objective. For seed solutions, pruning maximizes density.
For branch-and-bound incumbents, pruning maximizes the current parametric
objective. The implementation maintains an internal-degree cache so each removal
step can identify the least useful removable node efficiently. If connectivity
is being enforced, a small BFS-style check prevents removals that disconnect the
selected support from the query.

This pruning step is important because the constraint is $|S|\geq k$. The LP or
integer search may return a set larger than needed; pruning tightens the
incumbent and often reduces the returned community to a compact core.

\subsection{BFS Baseline}
\label{sec:bfs}

The \texttt{--bfs} variant performs topology-only frontier expansion from
$q$. Layered breadth-first: each dequeued node is oracle-queried, and its
predecessors and successors enter the candidate pool. A depth parameter
$d$ bounds the number of explored layers. There is no density objective,
no size constraint, and no $\kappa$-edge-connectivity enforcement, so the
returned set is the pool collected during expansion. The cost is one
oracle query per visited node, with no flow or LP work. Output size grows
with $d$; per-$k$ comparison against branch-and-price requires the
grow-to-$k$ post-processing of Section~\ref{sec:grow-to-k}.

\subsection{Average-Degree Baseline}
\label{sec:avgdeg}

The \texttt{--avgdeg} variant is an exact baseline for the query-anchored maximum
average-degree problem. Given the query node $q$, it returns
$$
  S^{\star} = \argmax_{S \ni q,\, S \subseteq V_{\mathrm{loc}}}\;
  \frac{|E(S)|}{|S|},
$$
where $E(S)$ counts directed edges with both endpoints in $S$ and
$V_{\mathrm{loc}}$ is the BFS-bounded local subgraph around $q$. The procedure
follows the Goldberg max-density flow construction \cite{goldberg1984} with three
adaptations.

\paragraph{Local oracle subgraph.}
Goldberg's algorithm operates on a fully materialised graph. Here the graph is
behind an oracle, so we first BFS from $q$ up to a depth bound (the
implementation default is unbounded). Each oracle query returns the node's
predecessors and successors; successors are cached during BFS and replayed in
one sweep to materialise the induced directed edge list. The output is a dense
indexing $(V_{\mathrm{loc}}, E_{\mathrm{loc}})$ with the query mapped to a known
index. Every subsequent step works on this local snapshot, not on the full
graph.

\paragraph{Parametric flow network.}
The Dinkelbach transform reduces the ratio objective to a parametric linear
form: for fixed $\lambda$,
$$
  \max_{S \ni q}\; |E(S)| - \lambda |S|.
$$
For each candidate $\lambda$, the solver builds a Goldberg-style $s$--$t$
network:
\begin{itemize}[leftmargin=1.1cm]
  \item a source $s$ and sink $t$;
  \item for each directed edge $e = (u,v) \in E_{\mathrm{loc}}$, an
  \emph{edge-item} node $n_e$ with $s \to n_e$ of capacity
  $\mathrm{SCALE} = 10^{6}$, and $n_e \to u$, $n_e \to v$ each of capacity
  $\infty$;
  \item for each node $i$, an arc $i \to t$ of capacity
  $\lfloor \lambda \cdot \mathrm{SCALE} \rfloor$;
  \item an anchor $s \to q$ of capacity $\infty$, which keeps $q$ on the source
  side of every minimum cut.
\end{itemize}
A push--relabel max-flow on this network and the source-side residual reachable
set give the optimal $S_\lambda$ for the parametric objective. The integer
scaling by $\mathrm{SCALE}$ keeps node capacities exact and bounds rounding
error by $1/\mathrm{SCALE}$.

\paragraph{Dinkelbach bisection.}
The solver bisects $\lambda$ on $[0, |E_{\mathrm{loc}}|]$ for up to one hundred
iterations or until the bracket falls below $10^{-9}$. After each probe it
records the achieved density $|E(S_\lambda)| / |S_\lambda|$ and keeps the
argmax across \emph{all} probed $\lambda$ rather than only the converged value.
A final fixup probes the bracket endpoints whenever they were not already the
last tested midpoint.

\paragraph{Differences from textbook Goldberg.}
Four substantive deviations from the classical construction matter for how the
baseline behaves on this problem.
\begin{enumerate}[leftmargin=1.1cm]
  \item \emph{Query anchor.} Goldberg returns the globally densest subset.
  Here the $s \to q$ arc forces $q$ into every cut, so the answer is the
  query-anchored densest subset. The two coincide only when the global optimum
  already contains $q$.
  \item \emph{Local subgraph.} Goldberg assumes a known input graph. The oracle
  setting forces a BFS-bounded local view, so the baseline maximises density on
  $V_{\mathrm{loc}}$, not on $V$. The depth bound is a hyperparameter, fixed in
  this work to ``full reachable from $q$''.
  \item \emph{Directed edges, undirected flow contribution.} Each directed edge
  $(u,v)$ produces a single edge-item connected to both endpoints with
  $\infty$ capacity, so a reciprocal pair $(u,v)$ and $(v,u)$ contributes two
  edge-items. The density numerator therefore counts directed edges. Goldberg's
  standard form takes undirected $\{u,v\}$ once.
  \item \emph{Density definition.} The avgdeg baseline reports
  $|E(S)| / |S|$ (directed edges over nodes), whereas the branch-and-price
  solver targets $|E(S)| / (|S|(|S|-1))$ (directed edge density). The two
  objectives are not monotone in each other: the avgdeg optimum on a graph need
  not be dense by the branch-and-price metric, and vice versa.
\end{enumerate}

The remaining mechanical differences are minor: a bisection bracket of
$[0, |E_{\mathrm{loc}}|]$ instead of $[0, \max\deg / 2]$, integer scaling
instead of exact rational capacities, and a consider-best loop instead of
Goldberg's interval-feasibility test. None of these change which subset is
optimal; they only affect numerical precision and termination.

\subsection{Grow-to-\texorpdfstring{$k$}{k} Post-Processing}
\label{sec:grow-to-k}

Branch-and-price enforces $|S| \ge k$ directly through the size constraint
$\sum_{u \in \Vactive} x_u \ge k$. BFS and AvgDeg do not have a size knob,
so a separate post-processing pass equalises their output cardinality with
branch-and-price for per-$k$ comparison. Given a seed $S$ and a candidate
pool $P \supseteq S$, the pass greedily adds the pool node with the
largest internal degree into the current set:
\[
  v^\star = \argmax_{v \in P \setminus S}\; \bigl|N(v) \cap S\bigr|,
\]
and repeats until $|S| = k$ or $P$ is exhausted. BFS uses $S = \{q\}$ with
depth $d$ acting as a minimum, and growth continues until the pool itself
reaches $k$. AvgDeg uses its Goldberg optimum $S^\star$ as the seed and
fires only when $|S^\star| < k$. The pass is a greedy heuristic; the
densest-at-least-$k$ problem is NP-hard, so no approximation guarantee is
claimed.

\subsection{End-to-End System}

The project is not only an algorithm implementation. The C++ solver is exposed
through a command-line interface with several modes: branch-and-price
(\texttt{--bp}), BFS (\texttt{--bfs}), exact average-degree baseline
(\texttt{--avgdeg}), and connected local baselines for synthetic experiments.
The FastAPI backend launches solver processes, streams log lines and structured
results, and supports stopping a running session. The React frontend exposes the
default branch-and-price path and lets a user set the query paper, $k$, time
limits, and column-generation parameters. The experimental CLI exposes the other
solver variants. The interface visualizes selected core papers and frontier
papers and shows a telemetry panel containing convergence and timing
information.

This interface is useful for research because dense-community optimization is
not a black-box recommendation. Users can inspect which papers were selected,
how many API queries were made, whether the solver is still improving, and where
time is being spent. The same telemetry also supports debugging and experiment
reproducibility.

\section{Experiments}

Evaluation has three perspectives that share per-query records. The first
is correctness on enumerable random graphs against brute force
(Section~\ref{sec:bf}). The second is intrinsic cluster quality and
solver cost on the full Cora-ML val+test pool
(Sections~\ref{sec:dataset}--\ref{sec:cost}). The third is downstream
node classification, tuned on val and evaluated on test plus a hard
test subset where the one-hop baseline already fails
(Sections~\ref{sec:classify-protocol}--\ref{sec:hard}). Cluster quality
is the headline perspective; classification is a downstream readout of
the retrieved geometry.

\subsection{Brute-Force Optimality on Small Random Graphs}
\label{sec:bf}

Before reporting on the citation graph, the branch-and-price solver is
checked against a brute-force oracle on synthetic graphs small enough to
enumerate every subset containing the query. The fixture is
Erd\H{o}s--R\'enyi $G(n, p)$ at $n = 20$, $p \in \{0.25, 0.50, 0.75\}$,
with seeds $0..4$ per $p$, giving $15$ graphs. The query node $q$ is the
highest-degree vertex of each graph.

Brute force enumerates every $S \ni q$ with $|S| \ge k$ for
$k \in \{3, 4, 5\}$. For each $k$ it records the maximum-density $S$ at
every connectivity threshold $\kappa \in \{0, 1, 2\}$. Cells with no
feasible $S$ at the requested $\kappa$ are marked infeasible and skipped
at the solver stage so the run stays finite. The C++ solver runs across
the same $(k, \kappa)$ grid per graph with no time, node, gap, or
Dinkelbach limit. A run is counted as a primary match when its objective
lies within $10^{-9}$ of the brute-force optimum. \texttt{optima.csv} and
\texttt{solver\_runs.csv} carry a \texttt{code\_hash} of the brute-force
source and a \texttt{solver\_build\_id} baked into the binary at
\texttt{cmake} time, so a stale brute force or stale solver binary is
caught at aggregation.

The result on the full $\{3, 4, 5\} \times \{0, 1, 2\}$ grid: AvgDeg
matches the unconstrained avg-degree optimum on $15/15$ graphs. BP
returns $S$ satisfying $|S| \ge k$ and rooted $\kappa$-edge-connectivity
on every feasible $(k, \kappa, p, \text{seed})$ cell, and matches the
brute-force edge density on every feasible cell. Zero solver errors
across the grid.

\subsection{Dataset and Split}
\label{sec:dataset}

The lead dataset is Cora-ML from PyTorch Geometric's \texttt{CitationFull}
collection with directed edges preserved: $2{,}995$ nodes, $8{,}416$
directed edges, $7$ paper labels. The query pool restricts to pure-source
papers (papers that cite others but are not themselves cited in the
exported graph) with out-degree at least $2$. With random seed $42$, the
pool splits evenly into $402$ val queries and $402$ test queries. The
remaining $2{,}191$ nodes form the train pool, which supplies labels for
the prediction rule and is never queried by any method. A narrower
\emph{hard subset} contains the $59$ test queries where the depth-$1$ BFS
majority vote disagrees with the true label; the subset is hash-pinned in
\texttt{split\_meta.json} so the same nodes are evaluated across runs.
CiteSeer, DBLP, and Cora follow the same protocol and are processed in
ascending node-count order.

All Cora-ML runs use \texttt{max\_in\_edges = 0}: every method expands
through outgoing citation edges only, matching the inductive split where
test papers cite older training papers. Enabling incoming edges would
change the retrieval problem and the label-propagation neighbourhood.

\subsection{Methods Compared}
\label{sec:methods}

Three solver variants are compared, in order of complexity:
\begin{itemize}[leftmargin=1.1cm]
  \item BFS at depth $1$ (Section~\ref{sec:bfs});
  \item AvgDeg, the exact average-degree baseline of
  Section~\ref{sec:avgdeg};
  \item BP, the branch-and-price solver, at $\kappa \in \{0, 1, 2\}$.
\end{itemize}
All three sweep $k \in \{3, 4, 5\}$. BFS and AvgDeg outputs are grown to
$|S| \ge k$ where feasible by the post-processing of
Section~\ref{sec:grow-to-k}, so the three methods return comparable
cardinality at each $k$. BP runs with soft time cap $60$ s and hard time
cap $300$ s; no explicit node, gap, or Dinkelbach limit.

\subsection{Cluster Quality on val + test}
\label{sec:cluster}

Cluster quality is the headline perspective.
Table~\ref{tab:cluster-quality} reports mean $\pm$ standard deviation
over the $804$ val+test queries on a single seed. Five intrinsic
statistics per query: subgraph size $|S|$, algebraic connectivity
$\lambda_2$ of the normalised Laplacian of the induced undirected
subgraph, normalised internal min-cut (ncut), average degree
$2|E(S)|/|S|$, and directed edge density $|E(S)| / (|S|(|S|-1))$.
$\lambda_2$ is computed on the dense Laplacian.

\begin{table}[t]
\centering
\footnotesize
\caption{Cluster quality on the Cora-ML val + test pool ($n = 804$
queries), mean $\pm$ std, single seed. \underline{Section best} (max
per column within each $k$, $|S|$ excluded); \textbf{overall best}
(max across $k$).}
\label{tab:cluster-quality}
\begin{tabular}{lccccc}
\toprule
Method & $|S|$ & $\lambda_2$ & ncut & avg.~deg & density \\
\midrule
\multicolumn{6}{l}{\textit{$k = 3$}} \\
BFS $d{=}1$       & $4.87 \pm 2.75$    & $0.96 \pm 0.30$ & $1.23 \pm 0.15$ & $1.17 \pm 0.50$ & $0.35 \pm 0.11$ \\
AvgDeg            & $21.20 \pm 15.20$  & $0.32 \pm 0.38$ & $0.70 \pm 0.53$ & \underline{\textbf{$3.02 \pm 1.55$}} & $0.21 \pm 0.12$ \\
BP $\kappa{=}0$   & $3.34 \pm 0.71$    & $1.23 \pm 0.34$ & $1.36 \pm 0.26$ & $1.20 \pm 0.42$ & \underline{\textbf{$0.51 \pm 0.10$}} \\
BP $\kappa{=}1$   & $3.32 \pm 0.67$    & $1.25 \pm 0.28$ & $1.39 \pm 0.13$ & $1.19 \pm 0.41$ & \underline{$0.51 \pm 0.10$} \\
BP $\kappa{=}2$   & $3.28 \pm 0.72$    & \underline{\textbf{$1.27 \pm 0.27$}} & \underline{\textbf{$1.41 \pm 0.11$}} & $1.10 \pm 0.44$ & $0.48 \pm 0.12$ \\
\midrule
\multicolumn{6}{l}{\textit{$k = 4$}} \\
BFS $d{=}1$       & $5.20 \pm 2.56$    & $0.78 \pm 0.25$ & $1.17 \pm 0.09$ & $1.22 \pm 0.47$ & $0.32 \pm 0.10$ \\
AvgDeg            & $21.21 \pm 15.18$  & $0.31 \pm 0.35$ & $0.69 \pm 0.53$ & \underline{$3.02 \pm 1.55$} & $0.21 \pm 0.11$ \\
BP $\kappa{=}0$   & $4.35 \pm 0.78$    & $0.75 \pm 0.48$ & $0.94 \pm 0.53$ & $1.59 \pm 0.48$ & \underline{$0.47 \pm 0.09$} \\
BP $\kappa{=}1$   & $4.23 \pm 0.76$    & $0.92 \pm 0.28$ & $1.21 \pm 0.10$ & $1.47 \pm 0.42$ & $0.46 \pm 0.09$ \\
BP $\kappa{=}2$   & $4.19 \pm 0.72$    & \underline{$0.94 \pm 0.28$} & \underline{$1.24 \pm 0.08$} & $1.34 \pm 0.46$ & $0.42 \pm 0.12$ \\
\midrule
\multicolumn{6}{l}{\textit{$k = 5$}} \\
BFS $d{=}1$       & $5.72 \pm 2.36$    & $0.66 \pm 0.24$ & $1.13 \pm 0.08$ & $1.29 \pm 0.46$ & $0.29 \pm 0.09$ \\
AvgDeg            & $21.25 \pm 15.14$  & $0.30 \pm 0.34$ & $0.69 \pm 0.52$ & \underline{$3.02 \pm 1.55$} & $0.20 \pm 0.11$ \\
BP $\kappa{=}0$   & $5.03 \pm 0.72$    & $0.50 \pm 0.44$ & $0.74 \pm 0.57$ & $1.76 \pm 0.51$ & \underline{$0.43 \pm 0.09$} \\
BP $\kappa{=}1$   & $5.14 \pm 1.03$    & \underline{$0.73 \pm 0.27$} & $1.14 \pm 0.08$ & $1.67 \pm 0.47$ & $0.41 \pm 0.09$ \\
BP $\kappa{=}2$   & $5.16 \pm 1.13$    & $0.72 \pm 0.28$ & \underline{$1.16 \pm 0.08$} & $1.52 \pm 0.54$ & $0.37 \pm 0.11$ \\
\bottomrule
\end{tabular}
\end{table}

At every $k$, BP $\kappa = 0$ takes the highest pair density and BP
$\kappa = 2$ takes the highest $\lambda_2$ and ncut, so BP wins internal
geometry across the grid. AvgDeg returns much larger sets
($|S| \approx 21$ regardless of $k$) with the highest average degree
($3.02$) by construction, but the spread across many more nodes drops
per-pair edge density to $0.20$--$0.21$. BFS depth-$1$ output size grows
from $4.9$ at $k = 3$ to $5.7$ at $k = 5$ with edge density between
$0.29$ and $0.35$, dominated on every internal axis.

\subsection{Cost on val + test}
\label{sec:cost}

Table~\ref{tab:cost} reports per-query wall time, oracle queries,
branch-and-bound nodes, LP solves, separated cuts, and LP time on the
same $804$ queries.

\begin{table}[t]
\centering
\footnotesize
\caption{Solver cost on the Cora-ML val + test pool ($n = 804$),
mean $\pm$ std. BP soft cap $60$ s, hard cap $300$ s.}
\label{tab:cost}
\begin{tabular}{lcccccc}
\toprule
Method & wall (s) & queries & BB nodes & LP solves & cuts & $t_{\text{LP}}$ (s) \\
\midrule
\multicolumn{7}{l}{\textit{$k = 3$}} \\
BFS $d{=}1$       & $0.03 \pm 0.01$  & $5 \pm 3$     & --                & --                  & --                  & --              \\
AvgDeg            & $0.08 \pm 0.05$  & $202 \pm 210$ & --                & --                  & --                  & --              \\
BP $\kappa{=}0$   & $1.54 \pm 15.19$ & $179 \pm 201$ & $121 \pm 837$     & $131 \pm 861$       & $87 \pm 315$        & $0.93 \pm 11.50$ \\
BP $\kappa{=}1$   & $1.47 \pm 13.71$ & $179 \pm 201$ & $121 \pm 772$     & $132 \pm 796$       & $90 \pm 308$        & $0.87 \pm 10.4$ \\
BP $\kappa{=}2$   & $10.3 \pm 53.1$  & $180 \pm 202$ & $81 \pm 429$      & $1106 \pm 5876$     & $1106 \pm 5868$     & $6.26 \pm 36.3$ \\
\midrule
\multicolumn{7}{l}{\textit{$k = 4$}} \\
BFS $d{=}1$       & $0.03 \pm 0.01$  & $5 \pm 3$     & --                & --                  & --                  & --              \\
AvgDeg            & $0.08 \pm 0.05$  & $202 \pm 210$ & --                & --                  & --                  & --              \\
BP $\kappa{=}0$   & $9.85 \pm 38.12$ & $200 \pm 210$ & $833 \pm 2025$    & $881 \pm 2079$      & $492 \pm 677$       & $9.21 \pm 36.31$ \\
BP $\kappa{=}1$   & $28.2 \pm 79.7$  & $200 \pm 210$ & $1549 \pm 3593$   & $1618 \pm 3696$     & $756 \pm 1348$      & $26.9 \pm 76.8$ \\
BP $\kappa{=}2$   & $47.9 \pm 103.3$ & $200 \pm 210$ & $1363 \pm 3061$   & $3794 \pm 9691$     & $3300 \pm 9483$     & $44.1 \pm 95.6$ \\
\midrule
\multicolumn{7}{l}{\textit{$k = 5$}} \\
BFS $d{=}1$       & $0.04 \pm 0.01$  & $6 \pm 2$     & --                & --                  & --                  & --              \\
AvgDeg            & $0.08 \pm 0.05$  & $202 \pm 210$ & --                & --                  & --                  & --              \\
BP $\kappa{=}0$   & $66.7 \pm 110.0$ & $202 \pm 210$ & $2967 \pm 4275$   & $3079 \pm 4392$     & $1291 \pm 1477$     & $63.2 \pm 105.6$ \\
BP $\kappa{=}1$   & $92.4 \pm 130.6$ & $202 \pm 210$ & $3799 \pm 5311$   & $3961 \pm 5485$     & $1992 \pm 2825$     & $88.4 \pm 126.3$ \\
BP $\kappa{=}2$   & $96.0 \pm 131.8$ & $202 \pm 210$ & $3197 \pm 4493$   & $4409 \pm 7318$     & $3326 \pm 6744$     & $90.8 \pm 126.7$ \\
\bottomrule
\end{tabular}
\end{table}

BFS is sub-second ($0.03$--$0.04$ s) with at most $6$ oracle queries.
AvgDeg costs $0.08$ s and $\approx 200$ queries; the query budget
reflects the full BFS-bounded local subgraph used by Goldberg's flow
construction. BP cost is heterogeneous: medians stay small but the mean
grows with $k$ and $\kappa$, from $1.5$ s at $k = 3, \kappa = 0$ to
$96$ s at $k = 5, \kappa = 2$. Standard deviations on wall time are
dominated by the tail of $\kappa = 2$ runs hitting the $300$ s hard cap.
$\kappa = 2$ also drives BQP separation up an order of magnitude
(cuts $\approx 1100$--$3300$ per query), and LP time tracks total wall
time closely.

\subsection{Classification Protocol}
\label{sec:classify-protocol}

For each val or test query $q$, the chosen method returns a community
$S \ni q$. The classifier predicts $q$'s label by uniform majority vote
over the training-labelled members of $S$, with ties broken by ascending
label value for reproducibility. The query itself and any val or test
node inside $S$ are excluded from the vote. If $S$ contains no
training-labelled member, the pipeline falls back to a concentric BFS
from $q$ that forbids val and test nodes as stepping stones until it
reaches a training node; if no training label is reachable, it returns
the global training-majority label. At the tuned hyperparameters, the
fallback fires on a negligible fraction of queries.

The retrieval set is therefore evaluated as a neighbourhood for label
propagation, not as input to a trained model. Differences in
classification metrics come from the selected community geometry, not
from learned features or model capacity.

\subsection{Tuning on val}
\label{sec:tune-val}

Table~\ref{tab:tune-val} reports macro precision, recall, and F1 on the
$402$ val queries across the $(k, \kappa)$ grid. The per-family val-best
setting is BFS depth-$1$ at $k = 4$, AvgDeg at $k = 4$, and BP at
$k = 4, \kappa = 2$. BFS depth-$1$ at $k = 4$ tops the table overall at
$F1 = 0.838 \pm 0.020$.

\begin{table}[t]
\centering
\footnotesize
\caption{Tuning on the Cora-ML val split ($n = 402$). Mean $\pm$
bootstrap std. \underline{Section best} (max per column within $k$);
\textbf{overall best} (max across $k$).}
\label{tab:tune-val}
\begin{tabular}{llll}
\toprule
Method                     & P     & R     & F1 \\
\midrule
\multicolumn{4}{l}{\textit{$k = 3$}} \\
BFS $d = 1$                & \underline{$0.836 \pm 0.019$} & \underline{$0.832 \pm 0.020$} & \underline{$0.826 \pm 0.021$} \\
AvgDeg                     & $0.714 \pm 0.026$ & $0.540 \pm 0.024$ & $0.558 \pm 0.026$ \\
BP $\kappa = 0$            & $0.797 \pm 0.021$ & $0.787 \pm 0.022$ & $0.782 \pm 0.022$ \\
BP $\kappa = 1$            & $0.799 \pm 0.020$ & $0.790 \pm 0.021$ & $0.785 \pm 0.021$ \\
BP $\kappa = 2$            & $0.801 \pm 0.020$ & $0.794 \pm 0.021$ & $0.785 \pm 0.022$ \\
\midrule
\multicolumn{4}{l}{\textit{$k = 4$}} \\
BFS $d = 1$                & \underline{\textbf{$0.844 \pm 0.019$}} & \underline{\textbf{$0.843 \pm 0.019$}} & \underline{\textbf{$0.838 \pm 0.020$}} \\
AvgDeg                     & $0.717 \pm 0.026$ & $0.544 \pm 0.024$ & $0.563 \pm 0.025$ \\
BP $\kappa = 0$            & $0.762 \pm 0.021$ & $0.683 \pm 0.025$ & $0.694 \pm 0.025$ \\
BP $\kappa = 1$            & $0.808 \pm 0.020$ & $0.789 \pm 0.021$ & $0.789 \pm 0.022$ \\
BP $\kappa = 2$            & $0.827 \pm 0.020$ & $0.810 \pm 0.021$ & $0.811 \pm 0.021$ \\
\midrule
\multicolumn{4}{l}{\textit{$k = 5$}} \\
BFS $d = 1$                & \underline{$0.814 \pm 0.021$} & \underline{$0.816 \pm 0.021$} & \underline{$0.810 \pm 0.021$} \\
AvgDeg                     & $0.717 \pm 0.026$ & $0.544 \pm 0.024$ & $0.563 \pm 0.025$ \\
BP $\kappa = 0$            & $0.725 \pm 0.023$ & $0.616 \pm 0.026$ & $0.628 \pm 0.026$ \\
BP $\kappa = 1$            & $0.772 \pm 0.024$ & $0.744 \pm 0.023$ & $0.741 \pm 0.024$ \\
BP $\kappa = 2$            & $0.795 \pm 0.022$ & $0.784 \pm 0.021$ & $0.780 \pm 0.022$ \\
\bottomrule
\end{tabular}
\end{table}

The temporal split favours BFS: pure-source test papers cite same-class
training papers, so a one-hop majority vote already retrieves the
homophilous signal. Within BP, $\kappa = 2$ dominates lower $\kappa$ at
every $k$ on the classification metric, even though density at
$\kappa = 2$ is below density at $\kappa = 0$. Connectivity therefore
helps the downstream vote without helping pair density.

\subsection{Test Classification}
\label{sec:test}

Table~\ref{tab:test-classification} applies the per-family val-best
setting to the $402$ Cora-ML test queries.

\begin{table}[t]
\centering
\footnotesize
\caption{Test classification on Cora-ML ($n = 402$). Per-family val-best
parameters carried over from Table~\ref{tab:tune-val}. Mean $\pm$
bootstrap std. \textbf{Bold} = max per column.}
\label{tab:test-classification}
\begin{tabular}{lllll}
\toprule
Method & val-best params      & P     & R     & F1 \\
\midrule
BFS    & $d = 1, k = 4$       & \textbf{$0.871 \pm 0.016$} & \textbf{$0.853 \pm 0.020$} & \textbf{$0.856 \pm 0.019$} \\
AvgDeg & $k = 4$              & $0.772 \pm 0.022$ & $0.579 \pm 0.024$ & $0.603 \pm 0.024$ \\
BP     & $k = 4, \kappa = 2$  & $0.818 \pm 0.021$ & $0.815 \pm 0.021$ & $0.811 \pm 0.021$ \\
\bottomrule
\end{tabular}
\end{table}

The val ranking carries over to test: BFS $0.856$, BP $0.811$, AvgDeg
$0.603$. AvgDeg's gap reflects its much larger communities diluting the
majority vote, not absent label information. BP closes most of the F1
gap to BFS while returning the densest compact community in
Table~\ref{tab:cluster-quality}.

\subsection{BFS-Hard Test Subset}
\label{sec:hard}

The hard subset contains the $59$ test queries where the depth-$1$ BFS
majority vote already disagrees with the true label.
Table~\ref{tab:hard-classification} applies the same per-family val-best
settings to this subset.

\begin{table}[t]
\centering
\footnotesize
\caption{Classification on the BFS-hard Cora-ML test subset ($n = 59$).
Mean $\pm$ bootstrap std. \textbf{Bold} = max per column.}
\label{tab:hard-classification}
\begin{tabular}{lllll}
\toprule
Method & val-best params      & P     & R     & F1 \\
\midrule
BFS    & $d = 1, k = 4$       & $0.183 \pm 0.061$ & $0.097 \pm 0.049$ & $0.080 \pm 0.042$ \\
AvgDeg & $k = 4$              & \textbf{$0.269 \pm 0.084$} & $0.251 \pm 0.079$ & $0.146 \pm 0.056$ \\
BP     & $k = 4, \kappa = 2$  & $0.174 \pm 0.053$ & \textbf{$0.282 \pm 0.090$} & \textbf{$0.190 \pm 0.056$} \\
\bottomrule
\end{tabular}
\end{table}

BFS scores near zero by construction: the subset is defined by depth-$1$
votes that are already wrong, and the grow-to-$k$ pass adds only
neighbours of those same wrong nodes. AvgDeg leads precision at $0.269$
on the back of much larger communities. BP leads recall ($0.282$) and F1
($0.190$) with the $k = 4, \kappa = 2$ community. Absolute F1 remains
below $0.2$ across the board: the subset is the difficult tail where the
cheapest baseline already fails.

\section{Conclusion}

\subsection{Discussion}

\paragraph{BP wins internal geometry at every $k$.}
On the $804$-query val + test pool, BP at $\kappa = 0$ returns the
densest community and BP at $\kappa = 2$ returns the highest $\lambda_2$
and ncut at every $k \in \{3, 4, 5\}$. The cluster-quality ranking is
stable across the grid: BP leads on pair density and algebraic
connectivity, BFS sits in the middle, AvgDeg trails on per-pair density.
The dense, compact, query-anchored community is the contribution;
classification metrics are a downstream readout of that geometry.

\paragraph{AvgDeg trades pair density for size.}
AvgDeg optimises average degree rather than pair density, and the two
objectives are not monotone in each other. The solver returns
$|S| \approx 21$ with average degree $3.02$ regardless of $k$, the
largest and best-connected set by average degree in the comparison. The
spread drops per-pair edge density to $0.20$--$0.21$ and the large set
dilutes the majority vote at classification time.

\paragraph{BFS depth-$1$ is the homophily ceiling for the easy split.}
Cora-ML test queries are pure-source papers whose direct citations land
on training papers. A one-hop majority vote matches that homophily
structure at the minimum oracle cost. Any retrieval method that deviates
from the one-hop frontier competes against a baseline already aligned
with the split.

\paragraph{The hard subset reorders the classifiers.}
On the $59$ test queries where the depth-$1$ vote is wrong by
construction, BFS scores $F1 = 0.080$, AvgDeg $0.146$, and BP $0.190$.
BP narrowly leads on the regime where the cheapest baseline already
fails. The cluster-quality ranking does not flip on the hard subset, so
the classification reorder is a regime-specific readout, not a geometry
change.

\paragraph{The role of $\kappa$.}
The classification-best BP setting uses $k = 4, \kappa = 2$. Within BP,
$\kappa = 2$ dominates $\kappa \in \{0, 1\}$ on F1 at every $k$, even
though density at $\kappa = 2$ is slightly below density at $\kappa = 0$.
$\kappa$ is a tuning parameter that trades pair density for structural
coherence; it is not universally beneficial.

\paragraph{Cost grows steeply with $k$ and $\kappa$.}
BFS is sub-second. AvgDeg costs $0.08$ s and $\approx 200$ queries. BP
mean wall time grows from $1.5$ s at $k = 3, \kappa = 0$ to $96$ s at
$k = 5, \kappa = 2$ under the $300$ s hard cap; $\kappa = 2$ also drives
BQP separation up an order of magnitude. The live OpenAlex telemetry
exposes that cost-quality tradeoff at interactive latency.

\subsection{Limitations and Future Work}

Reported numbers are Cora-ML only at this draft. The same pipeline runs
on CiteSeer, DBLP, and Cora; per-dataset writeups are produced in
ascending node-count order. The temporal split makes one-hop BFS
unusually strong: test queries are pure-source papers whose direct
citations sit in the train set, so the easy regime is wide. The hard
subset is a conditional diagnostic on test queries where the depth-$1$
BFS majority vote is wrong, so the BFS solver scores near zero on it by
construction.

The solver depends on Gurobi. The OpenAlex live mode inherits OpenAlex
coverage and entity-disambiguation limits. The lazy expansion procedure
optimises over the discovered graph region; a formal global optimality
claim requires assumptions about oracle completeness, so the
implementation is best viewed as a practical exact method over an
adaptively discovered region. Solver limits (time, branch-and-bound
nodes, Dinkelbach and column-generation iterations) return possibly
suboptimal incumbents on the tail of the run distribution.

Future work:
\begin{enumerate}[leftmargin=1.2cm]
  \item complete the CiteSeer, DBLP, and Cora runs and update the
  per-dataset tables;
  \item tighter LP relaxation and inverse-out-degree edge weights to
  attenuate hub effects;
  \item warm-start with a densest-subgraph incumbent to reduce
  Dinkelbach iterations;
  \item incorporate semantic information (abstract embeddings, author or
  venue affinity) into the objective;
  \item user studies for interactive paper discovery where compactness
  and interpretability matter more than classification accuracy.
\end{enumerate}

\subsection{Summary}

This paper presents a query-anchored densest community search system for
directed citation graphs. The method uses Dinkelbach fractional
programming, lazy branch-and-price, frontier pricing, branch-and-bound,
Boolean quadratic polytope triangle cuts, incumbent pruning, and optional
rooted $\kappa$-edge-connectivity cuts. On the Cora-ML val + test pool,
BP returns the densest and most algebraically connected compact community
at every $k \in \{3, 4, 5\}$. BFS depth-$1$ wins macro F1 on the
$402$-query test split because of direct citation homophily; BP runs
second ($0.811$ versus $0.856$). On the $59$-node BFS-hard test subset,
BP leads macro F1 ($0.190$) ahead of AvgDeg ($0.146$) and BFS ($0.080$).
BP is therefore a dense-community retriever; classification F1 is a
downstream readout of the retrieved geometry.

\subsection*{Reproducibility}

The reported artifacts are generated by the repository scripts. The
environment used for Python verification is the \texttt{dcs} conda
environment. The grid is $k \in \{3, 4, 5\}$ and, for BP,
$\kappa \in \{0, 1, 2\}$; BP runs with soft cap $60$ s, hard cap
$300$ s, and no node, gap, or Dinkelbach limit.

\begin{verbatim}
./solver/build.sh
conda run -n dcs python scripts/classification/cluster_quality.py \
  --dataset Cora_ML --splits val test
conda run -n dcs python scripts/classification/tune_methods.py \
  --dataset Cora_ML --family all --optimize f1 \
  --k_grid 3 4 5 --kappa_grid 0 1 2 \
  --bp_time_limit 60 --bp_hard_time_limit 300 \
  --bp_cg_batch_frac 1.0 --bp_cg_min_batch 50 --bp_cg_max_batch 50 \
  --bp_node_limit -1 --bp_gap_tol 0.0 --bp_dinkelbach_iter -1
conda run -n dcs python scripts/classification/evaluate_fixed_methods.py \
  --dataset Cora_ML --split test \
  --settings_json exps/classification/Cora_ML/tune/best_settings_all.json
conda run -n dcs python scripts/classification/evaluate_fixed_methods.py \
  --dataset Cora_ML --split test --subset bfs_depth1_wrong \
  --settings_json exps/classification/Cora_ML/tune/best_settings_all.json
conda run -n dcs python -m unittest scripts.classification.test_solver_utils -v
\end{verbatim}

\begin{thebibliography}{10}
\sloppy

\bibitem{goldberg1984}
A.~V. Goldberg.
\newblock Finding a maximum density subgraph.
\newblock Technical Report UCB/CSD-84-171, EECS Department, University of
California, Berkeley, 1984.
\newblock \url{https://www2.eecs.berkeley.edu/Pubs/TechRpts/1984/5956.html}.

\bibitem{charikar2000}
M. Charikar.
\newblock Greedy approximation algorithms for finding dense components in a
graph.
\newblock In \emph{Approximation Algorithms for Combinatorial Optimization},
LNCS 1913, pages 84--95, 2000.
\newblock \url{https://dblp.org/rec/conf/approx/Charikar00.html}.

\bibitem{dinkelbach1967}
W. Dinkelbach.
\newblock On nonlinear fractional programming.
\newblock \emph{Management Science}, 13(7):492--498, 1967.
\newblock \url{https://doi.org/10.1287/mnsc.13.7.492}.

\bibitem{barnhart1998}
C. Barnhart, E.~L. Johnson, G.~L. Nemhauser, M.~W.~P. Savelsbergh, and
P.~H. Vance.
\newblock Branch-and-price: Column generation for solving huge integer
programs.
\newblock \emph{Operations Research}, 46(3):316--329, 1998.
\newblock \url{https://doi.org/10.1287/opre.46.3.316}.

\bibitem{veremyev2017}
A. Veremyev, O.~A. Prokopyev, E.~L. Pasiliao, and V. Boginski.
\newblock Exact and superpolynomial approximation algorithms for the densest
$k$-subgraph problem.
\newblock \emph{European Journal of Operational Research}, 262(3):894--903,
2017.
\newblock \url{https://doi.org/10.1016/j.ejor.2017.04.034}.

\bibitem{priem2022}
J. Priem, H. Piwowar, and R. Orr.
\newblock OpenAlex: A fully-open index of scholarly works, authors, venues,
institutions, and concepts.
\newblock arXiv:2205.01833, 2022.
\newblock \url{https://arxiv.org/abs/2205.01833}.

\bibitem{openalexworks}
OpenAlex technical documentation.
\newblock Works API entity documentation.
\newblock \url{https://docs.openalex.org/api-entities/works}.

\bibitem{fey2019}
M. Fey and J.~E. Lenssen.
\newblock Fast graph representation learning with PyTorch Geometric.
\newblock arXiv:1903.02428, 2019.
\newblock \url{https://arxiv.org/abs/1903.02428}.

\bibitem{pygcitationfull}
PyTorch Geometric documentation.
\newblock \texttt{torch\_geometric.datasets.CitationFull}.
\newblock \url{https://pytorch-geometric.readthedocs.io}.

\bibitem{bojchevski2017}
A. Bojchevski and S. Guennemann.
\newblock Deep Gaussian embedding of graphs: Unsupervised inductive learning
via ranking.
\newblock arXiv:1707.03815, 2017.
\newblock \url{https://arxiv.org/abs/1707.03815}.

\end{thebibliography}

\end{document}
